% !PRINTED 64
\weberpage{64}
%<<CONT
$\nu$ is replaced by $\nu + \nu'$, if $\nu'$ denotes the index of $n'$.
By this there follows from (3):
\[
\varepsilon^{-\beta \nu'} (\varepsilon^{\beta}, r) = \sum^{n} \varepsilon^{\beta \nu} r^{n n'},
\]
and if one therefore multiplies by $r^{n'}$ and takes the sum over $n'$:
\weq{4}{(\varepsilon^{-\beta}, r)(\varepsilon^{\beta}, r) = \sum^{n, n'} \varepsilon^{\beta \nu} r^{n'(n+1)}.}
Thus first of all the sum
\weq{5}{\sigma = \sum^{n'} r^{n'(n+1)}}
is to be determined. But the sum $\sigma$ always vanishes when
$n + 1$ is not divisible by $q^{\varkappa - 1}$ (by Bd.\ I, \S.\,133); for
then $r^{n+1}$ is a root of unity whose degree is a higher than
the first power of $q$. Hence $\sigma$ can be different from zero
only when $n$ has the form
\weq{6}{n \equiv -1 + t q^{\varkappa - 1} \pmod{m},}
and then we shall denote its value by $\sigma_t$. We obtain
all values of $n$ with respect to the modulus $m$ that are contained
in the form (6) if we let $t$ run through a complete system of residues
with respect to the modulus $q$, so for instance the values
\weq{7}{t = 0, 1, 2 \ldots q - 1}
may be set.
The sum $\sigma$ consists of $\varphi(m)$ terms. If $t = 0$, then
each of these terms becomes $= 1$ and it follows that
\weq{8}{\sigma_0 = \varphi(m) = q^{\varkappa} - q^{\varkappa - 1};}
for every other value of $t$, $r^{n+1}$ is a primitive $q$th root of unity
, and in $\sigma_t$ every such root of unity then occurs
$\varphi(m) : (q - 1) = q^{\varkappa - 1}$ times. But the sum of all primitive
$q$th roots of unity is equal to $-1$, and hence it follows for
$t = 1, 2 \ldots q - 1$ that
\weq{9}{\sigma_t = -q^{\varkappa - 1}.}
With this the sum $\sigma$ is determined.

But in order to derive from this the value of the sum in (4),
it is still necessary to ascertain the index $\nu$ of the numbers $n$ of the form (6).

For these numbers, however,
\[
g^{\nu} \equiv -1 \pmod{q^{\varkappa - 1}},
\]
and since by \gesperrt{Fermat's} theorem [\S.\,14, (11)]
\weq{10}{g^{1/2 \varphi(q^{\varkappa})} \equiv -1 \pmod{q^{\varkappa}}}
%CONT>>
